% SPDX-License-Identifier: GPL-3.0-or-later
% Copyright (C) 2026 Alessandro Greco
\chapter{Построение цепи: TAP, CREATE\_FAST, ntor и ntor-v3}\label{cap:09-circuiti}
\aggiornatoal{23 сентября 2026 г.}{коммит torspec \texttt{8e1f245}; tor 0.4.9.13; предложения 216, 332, 350}

\begin{obiettivi}
К концу этой главы вы сможете:
\begin{itemize}[nosep]
  \item описывать, как клиент строит цепь по одному переходу с CREATE2 и EXTEND2;
  \item объяснять, что должен гарантировать протокол согласования ключей цепи: аутентификацию реле, секретность и прямую секретность;
  \item пошагово разбирать согласование ключей CREATE\_FAST, ntor и ntor-v3 и понимать, почему TAP был выведен из эксплуатации;
  \item получать четыре ключа одного перехода (Df, Db, Kf, Kb) из согласованного ключевого материала;
  \item проверять реализацию по официальным тестовым векторам;
  \item строить реальные цепи к реле лабораторной сети (лабораторная работа~3).
\end{itemize}
\end{obiettivi}

Имея открытый канал (\cref{cap:07-canali}) и алфавит ячеек (\cref{cap:08-celle}), мы располагаем всем необходимым для построения центральной структуры Tor: \textbf{цепи}. В этой главе мы увидим, как клиент устанавливает с каждым реле маршрута общий секрет, неизвестный никому другому, и как превращает этот секрет в ключи, шифрующие трафик. Использование этих ключей для послойного шифрования рассматривается в \cref{cap:10-cifratura-relay}.

% ---------------------------------------------------------------------------
\section{По одному переходу}\label{sec:circ-telescopio}

Клиент строит цепь \emph{постепенно}, или \enquote{телескопически}~\cite{dingledine2004tor}:
\begin{enumerate}
  \item Выбирается маршрут, например входной узел $R_1$ (guard), промежуточный $R_2$ (middle) и выходной $R_3$ (exit) (\cref{cap:13-percorso}).
  \item По каналу к $R_1$ клиент отправляет ячейку \cella{CREATE2}, содержащую первую половину криптографического протокола согласования ключей. $R_1$ отвечает \cella{CREATED2}, после чего клиент и $R_1$ имеют общий ключевой материал.
  \item Клиент отправляет сообщение \texttt{EXTEND2}, зашифрованное для $R_1$, с просьбой расширить цепь до $R_2$. В сообщении находится первая половина нового протокола согласования ключей, предназначенного уже $R_2$.
  \item $R_1$ открывает или повторно использует канал к $R_2$ и отправляет ему ячейку \cella{CREATE2} с этими данными согласования ключей. Ответ \cella{CREATED2} возвращается к $R_1$, который пересылает его клиенту в сообщении \texttt{EXTENDED2}.
  \item Клиент завершает согласование ключей с $R_2$ и повторяет процедуру для $R_3$, на этот раз через $R_1$ и $R_2$.
\end{enumerate}

\begin{figure}[htbp]
\centering
% SPDX-License-Identifier: GPL-3.0-or-later
% Copyright (C) 2026 Alessandro Greco
%% Русское альтернативное описание EPUB.
\altfig{Телескопическое построение цепи Tor из трёх переходов: клиент, R1 (сторожевой), R2 (промежуточный), R3 (выходной) показаны отдельными колонками. Клиент обменивается CREATE2/CREATED2 с R1 и выводит ключи K1. Затем EXTEND2 идёт внутри существующей цепи: R1 обрабатывает запрос и отправляет CREATE2 к R2, а ответы CREATED2/EXTENDED2 возвращают данные для вывода K2. Для третьего перехода EXTEND2 проходит через сторожевой узел к R2; R2 обрабатывает его и посылает CREATE2 к R3, после чего клиент выводит K3. Каждое рукопожатие создаёт отдельные ключи между клиентом и целевым узлом; уже задействованные узлы пересылают сообщения для следующих переходов, обработав собственный слой. Пунктирные стрелки изображают логический путь через промежуточные узлы, а не прямое соединение. Между соседними узлами действуют отдельные TLS-каналы.}
% Телескопическое построение цепи.
\begin{tikzpicture}[font=\sffamily\small,>=Stealth,
  msg/.style={font=\ttfamily\scriptsize,fill=white,inner sep=1pt},
  indietro/.style={-{Stealth[length=2.5mm]},thick,torverde!60!black},
  circav/.style={cellaf,dashed},
  circind/.style={indietro,dashed}]
  \foreach \n/\x/\t in {c/0/Клиент,r1/3.4/$R_1$ сторожевой,r2/6.8/$R_2$ промежуточный,r3/10.2/$R_3$ выходной}
    \node[font=\sffamily\small] (\n) at (\x,0) {\t};
  \foreach \x in {0,3.4,6.8,10.2} \draw[torgrigio] (\x,-0.3) -- (\x,-5.5);
  % первый переход
  \draw[cellaf] (0,-0.9) -- node[msg,above]{CREATE2} (3.4,-0.9);
  \draw[indietro] (3.4,-1.5) -- node[msg,above]{CREATED2} (0,-1.5);
  \node[anchor=east,font=\scriptsize,text=torgrigio] at (-0.2,-1.2) {ключи $K_1$};
  % второй переход
  \draw[circav] (0,-2.5) -- node[msg,above]{EXTEND2} (3.4,-2.5);
  \draw[cellaf] (3.4,-2.5) -- node[msg,above]{CREATE2} (6.8,-2.5);
  \draw[indietro] (6.8,-3.1) -- node[msg,above]{CREATED2} (3.4,-3.1);
  \draw[circind] (3.4,-3.1) -- node[msg,above]{EXTENDED2} (0,-3.1);
  \node[anchor=east,font=\scriptsize,text=torgrigio] at (-0.2,-2.8) {ключи $K_2$};
  % третий переход
  \draw[circav] (0,-4.2) -- node[msg,above]{EXTEND2} (6.8,-4.2);
  \draw[cellaf] (6.8,-4.2) -- node[msg,above]{CREATE2} (10.2,-4.2);
  \draw[indietro] (10.2,-4.8) -- node[msg,above]{CREATED2} (6.8,-4.8);
  \draw[circind] (6.8,-4.8) -- node[msg,above]{EXTENDED2} (0,-4.8);
  \node[anchor=east,font=\scriptsize,text=torgrigio] at (-0.2,-4.5) {ключи $K_3$};
  % каналы
  \foreach \a/\b/\y in {0/3.4/-5.8,3.4/6.8/-5.8,6.8/10.2/-5.8}
    \draw[thick,torgrigio!60,{Bar}-{Bar}] (\a,\y) -- node[below,font=\scriptsize]{TLS-канал} (\b,\y);
  \node[font=\scriptsize,text=torgrigio,align=center] at (5.1,-6.6) {каждое рукопожатие — между клиентом и одним узлом};
\end{tikzpicture}

\caption{Телескопическое построение трёхпереходной цепи. Каждый протокол согласования ключей выполняется только между клиентом и одним реле; промежуточные реле лишь переносят сообщения. Пунктирными стрелками показаны сообщения, идущие внутри цепи и зашифрованные ключами уже установленных переходов.}
\label{fig:telescopio}
\end{figure}

У такого устройства есть фундаментальное следствие. Согласование ключей с $R_3$ проходит через $R_1$ и $R_2$, но его сообщения зашифрованы для открытого ключа $R_3$: ни $R_1$, ни $R_2$ не могут получить его общий секрет. Поскольку каждый переход использует отдельный секрет, ни одно реле не может прочитать то, чем клиент обменивается с другими. Входной узел $R_1$ знает клиента, но не знает, куда направляется трафик; выходной $R_3$ знает пункт назначения, но не знает отправителя.

\begin{storia}[Onion routing первого поколения]
В исходной системе Naval Research Laboratory клиент строил одну \enquote{луковицу}: сообщение, послойно зашифрованное открытыми ключами каждого реле. Каждое реле снимало слой, чтобы узнать следующий переход и собственный сеансовый ключ (\cref{cap:03-nascita-tor}). Реле, скомпрометированное в будущем, могло бы расшифровать ранее записанные «луковицы». Телескопическая цепь Tor с эфемерным обменом Диффи--Хеллмана для каждого перехода ввела \emph{прямую секретность}: после удаления сеансовых ключей даже владение долговременными ключами не позволяет расшифровать прошлый трафик~\cite{dingledine2004tor}.
\end{storia}

% ---------------------------------------------------------------------------
\section{Ячейки и сообщения}\label{sec:circ-celle}

\subsection{CREATE2 и CREATED2}

Формат \cella{CREATE2} намеренно универсален~\spec{tor-spec}{Create and Created cells}: тип протокола согласования ключей (HTYPE, 2 байта), длина (HLEN, 2 байта) и данные (HDATA). \cella{CREATED2} содержит только HLEN и HDATA, поскольку тип уже известен. Всё помещается в ячейку фиксированной длины, поэтому данные согласования ключей в \cella{CREATE2} не могут превышать $509 - 4 = 505$ байт.

\begin{center}
\begin{tabular}{cll}
\toprule
HTYPE & Протокол & Статус \\
\midrule
\texttt{0x0000} & TAP & устарел; удалён в tor 0.4.9.1-alpha \\
\texttt{0x0001} & --- & зарезервировано \\
\texttt{0x0002} & ntor & используется \\
\texttt{0x0003} & ntor-v3 & используется; предпочтителен \\
\bottomrule
\end{tabular}
\end{center}

Если согласование ключей прошло успешно, реле отвечает \cella{CREATED2}; при ошибке --- \cella{DESTROY}. CircID новой ячейки следует правилам \cref{sec:circid}: начиная с версии~4 протокола соединения сторона, открывшая канал, использует значения со старшим битом 1.

\subsection{EXTEND2 и EXTENDED2}\label{sec:extend2}

Сообщение RELAY \texttt{EXTEND2} (команда RELAY 14) помимо полей \cella{CREATE2} содержит список \emph{спецификаторов соединения (link specifiers)}, сообщающих реле, как достичь следующего перехода и как его идентифицировать~\spec{tor-spec}{Extend and Extended messages}:

\begin{center}
\begin{tabular}{cll}
\toprule
Тип & Содержимое & Длина \\
\midrule
\texttt{00} & IPv4-адрес и ORPort & 6 \\
\texttt{01} & IPv6-адрес и ORPort & 18 \\
\texttt{02} & устаревшая RSA-идентичность (SHA-1, 20 байт) & 20 \\
\texttt{03} & идентичность Ed25519 & 32 \\
\bottomrule
\end{tabular}
\end{center}

Реле, расширяющее цепь, должно проверить \emph{каждую} указанную идентичность и не должно расширять цепь, если достигнутое реле не докажет владение ими во время установления канала из \cref{cap:07-canali}. Так клиент уверен, что сообщение согласования ключей дошло именно до выбранного реле, хотя напрямую с $R_2$ не общается. Реле также отклоняет запрос на расширение цепи к самому себе или к реле, от которого запрос поступил.

Две детали показывают, насколько внимательно протокол относится к \emph{неразличимости}. Во-первых, спецификация рекомендует всегда отправлять спецификаторы в одном порядке (\texttt{00}, \texttt{02}, \texttt{03}, \texttt{01}), чтобы порядок не выдавал реализацию. Во-вторых, \texttt{EXTEND2} всегда передаётся внутри \cella{RELAY\_EARLY} (\cref{sec:comandi}); ограниченное число таких ячеек не позволяет строить произвольно длинные цепи.

\texttt{EXTENDED2} (команда RELAY 15) имеет в точности то же тело, что \cella{CREATED2}.

% ---------------------------------------------------------------------------
\section{Что должно обеспечивать согласование ключей}\label{sec:circ-requisiti}

Перед рассмотрением разных протоколов согласования ключей определим ожидаемые свойства. По завершении клиент и реле должны иметь общий секрет со следующими свойствами:
\begin{description}
  \item[Аутентификация реле] Клиент должен знать, что делит секрет именно с \emph{конкретным} реле и никем больше. Реле, напротив, не должно узнавать ничего о клиенте: согласование ключей аутентифицировано \emph{только в одном направлении}; иначе каждое реле узнавало бы идентичность клиента.
  \item[Секретность] Сетевой наблюдатель или предыдущее реле цепи, пересылающее сообщения, не должны уметь вывести секрет.
  \item[Прямая секретность] Даже тот, кто в будущем получит долговременные ключи реле, не должен иметь возможности восстановить прошлые сеансовые секреты.
  \item[Эффективность] Нагруженное реле выполняет тысячи согласований ключей в секунду, поэтому вычислительная стоимость имеет значение.
\end{description}
Каждый рассматриваемый протокол согласования ключей по-своему балансирует эти требования.

% ---------------------------------------------------------------------------
\section{TAP: первый протокол согласования ключей}\label{sec:tap}

Исходный протокол согласования ключей Tor TAP (\emph{Tor Authentication Protocol}) сочетал обычный обмен ключами Диффи--Хеллмана в 1024-битной модульной группе с шифрованием открытым 1024-битным RSA-ключом. Клиент шифровал свою половину Диффи--Хеллмана RSA \enquote{onion}-ключом реле; расшифровать её и ответить могло только выбранное реле. Затем ключи вырабатывались функцией KDF-TOR на основе SHA-1 (\cref{sec:circ-chiavi}).

TAP удовлетворял требованиям, но с постоянно уменьшающимся запасом безопасности: 1024 бит недостаточно и для RSA, и для модульного Диффи--Хеллмана, гибридная схема шифрования считалась хрупкой, а согласования ключей были вычислительно дорогими. С 2013 года TAP дополнил ntor, а начиная с tor 0.4.9.1-alpha реле полностью перестали его поддерживать (\proposta{350}). Один след остаётся: ради совместимости реле всё ещё публикуют старый RSA \enquote{onion}-ключ в дескрипторах. В лабораторной сети он виден в микродескрипторах как \texttt{onion-key}, рядом с реально используемым \texttt{ntor-onion-key}.

% ---------------------------------------------------------------------------
\section{CREATE\_FAST: когда открытый ключ не нужен}\label{sec:create-fast}

Есть один случай, когда согласование ключей с открытым ключом не нужно: \emph{первый} переход. Клиент только что открыл TLS-канал к этому реле и проверил его идентичность, поэтому TLS уже даёт зашифрованный аутентифицированный канал. \cella{CREATE\_FAST} использует это обстоятельство~\spec{tor-spec}{CREATE\_FAST/CREATED\_FAST cells}:
\begin{enumerate}
  \item Клиент отправляет 20 случайных байт $X$.
  \item Реле отвечает 20 случайными байтами $Y$ и 20 байтами $KH$.
  \item Обе стороны получают ключевой материал функцией KDF-TOR из $K_0 = X \| Y$. Первые 20 байт должны совпасть с $KH$, доказывая, что реле получило $X$.
\end{enumerate}

Безопасность \cella{CREATE\_FAST} целиком зависит от TLS-канала: тот, кто способен прочитать TLS-трафик, увидит и $X$, и $Y$. Поэтому его следует использовать только для одношаговых цепей. Типичный случай --- запускающийся клиент, которому нужно загрузить каталожную информацию с реле, onion key которого он ещё не знает. Исходный код tor говорит об этом явно: \file{should\_use\_create\_fast\_for\_circuit} (в \file{src/core/or/circuitbuild.c}) выбирает \cella{CREATE\_FAST} только при отсутствии пригодного onion key первого перехода.

\begin{attenzione}[CREATE\_FAST раскрывает первый переход]
Реле, получившее \cella{CREATE\_FAST}, достоверно знает, что является первым переходом цепи и, следовательно, что сторона на другом конце канала --- клиент. Спецификация предлагает реле использовать именно этот признак, чтобы отклонять одношаговые цепи, пытающиеся выйти в Интернет.
\end{attenzione}

% ---------------------------------------------------------------------------
\section{ntor}\label{sec:ntor}

Протокол ntor, предложенный Голдбергом, Стебилой и Устаоглу~\cite{goldberg2013ntor} и введённый в Tor с \proposta[Предложением]{216} (tor 0.2.4.8-alpha, 2013), заменяет RSA и модульный Диффи--Хеллман на \textbf{Curve25519}, быстрый обмен ключами на эллиптической кривой с ключами всего по 32 байта (\cref{cap:05-crittografia}).

\subsection{Компоненты}

Реле имеет ключ onion среднего срока действия на Curve25519: секретный $b$ и открытый $B = b\,G$, опубликованный в дескрипторе как \texttt{ntor-onion-key}. Клиент знает $B$ и идентичность реле (здесь 20-байтовый RSA-отпечаток \texttt{ID}). Через $\mathrm{EXP}(P, k)$ обозначим обмен Диффи--Хеллмана Curve25519 между точкой $P$ и скаляром $k$, а через $H(x, t)$ --- HMAC-SHA256 сообщения $x$ с ключом $t$. Ключи $t$ --- постоянные строки, полученные из идентификатора протокола \texttt{"ntor-curve25519-sha256-1"}.

\subsection{Протокол}

\begin{enumerate}
  \item Клиент генерирует эфемерную пару $(x, X)$ и отправляет в \cella{CREATE2}: \texttt{ID} (20 байт), $B$ (32 байта, определяющие использованный ключ onion) и $X$ (32 байта). Всего 84 байта.
  \item Реле проверяет \texttt{ID} и $B$, генерирует эфемерную пару $(y, Y)$ и вычисляет
  % В EPUB формула становится SVG; доступный альтернативный текст задаётся
  % в epub.mk4, поскольку tex4ht не применяет \altfig к display math.
  % Любое изменение формулы требует проверки соответствующего alt-текста.
  \[ \mathit{secret\_input} = \mathrm{EXP}(X, y) \,\|\, \mathrm{EXP}(X, b) \,\|\, \texttt{ID} \,\|\, B \,\|\, X \,\|\, Y \,\|\, \texttt{PROTOID}. \]
  Из этого выводится $\mathit{KEY\_SEED} = H(\mathit{secret\_input}, t_{\mathit{key}})$ и подтверждающее значение $\mathit{AUTH}$. Ответ содержит $Y$ (32 байта) и $\mathit{AUTH}$ (32 байта).
  \item Клиент вычисляет тот же $\mathit{secret\_input}$ со своими секретами: $\mathrm{EXP}(Y, x) \,\|\,\allowbreak \mathrm{EXP}(B, x) \,\|\, \ldots$ --- и проверяет совпадение $\mathit{AUTH}$.
\end{enumerate}

В основе согласования ключей лежат два обмена Диффи--Хеллмана:
\begin{itemize}
  \item $\mathrm{EXP}(X, y) = \mathrm{EXP}(Y, x)$ сочетает два \textbf{эфемерных} ключа и обеспечивает прямую секретность, поскольку $x$ и $y$ удаляются сразу после согласования ключей.
  \item $\mathrm{EXP}(X, b) = \mathrm{EXP}(B, x)$ включает ключ onion реле и обеспечивает \textbf{аутентификацию}: только сторона, знающая $b$, может вычислить значение и создать правильный $\mathit{AUTH}$.
\end{itemize}
Ни одного из этих обменов недостаточно по отдельности. Без первого тот, кто похитит $b$ в будущем, сможет восстановить всё; без второго посредник сможет выдать себя за реле. Клиент при этом не использует долговременный ключ: реле знает лишь, что делит секрет с \enquote{тем, кто отправил $X$}.

\begin{approfondimento}[Проверка точки на бесконечности]
Спецификация требует от обеих сторон проверять, что ни один обмен Диффи--Хеллмана не даёт \enquote{точку на бесконечности}, представленную здесь полностью нулевым результатом. Противник мог бы отправить точку малого порядка как $Y$, заставив $\mathrm{EXP}(Y, x)$ принять предсказуемое значение. Для Curve25519 этой простой проверки достаточно вместо полной проверки принадлежности $Y$ правильной группе. В лабораторной работе~3 это три строки функции \file{x25519\_dh}.
\end{approfondimento}

% ---------------------------------------------------------------------------
\section{ntor-v3}\label{sec:ntor3}

У ntor есть ограничение: кроме ключей он ничего не переносит. Со временем появилась необходимость \emph{согласовывать} параметры при установлении перехода, например управление перегрузкой (\cref{cap:11-stream}) или CGO (\cref{cap:10-cifratura-relay}). ntor-v3 (\proposta{332}) расширяет ntor \textbf{зашифрованными сообщениями} в обоих направлениях и одновременно меняет ряд криптографических решений.

\subsection{Что меняется}

\begin{itemize}
  \item Идентичность реле теперь --- его \textbf{Ed25519}-идентичность (32 байта), а не RSA-отпечаток.
  \item Хеш-функции относятся к семейству SHA-3: SHA3-256 для хеширования и MAC, SHAKE-256 для выработки ключей переменной длины. Перед каждым входом добавляется доменная метка с кодированной длиной (функция $\mathrm{ENCAP}$), поэтому значение для одной цели нельзя принять за значение для другой.
  \item Строка проверки \texttt{VER}, известная обеим сторонам, участвует в каждом вычислении: согласование ключей из одного контекста нельзя повторно использовать в другом. Для обычных цепей tor применяет \texttt{"circuit extend"} (константа \file{NTOR3\_CIRC\_VERIFICATION} в \file{src/core/crypto/onion\_crypto.c}).
  \item Клиент может передать сообщение $CM$, а реле --- сообщение $SM$; оба зашифрованы.
\end{itemize}

\subsection{Протокол}

\begin{enumerate}
  \item Клиент генерирует $(x, X)$ и сразу вычисляет $B_x = \mathrm{EXP}(B, x)$. Из $B_x$ и публичных параметров он получает два ключа \enquote{первой фазы}: один для шифрования $CM$ посредством AES-256-CTR, второй для его аутентификации MAC. Отправляются \texttt{ID} (32), $B$ (32), $X$ (32), зашифрованный $CM$ и MAC (32).
  \item Реле вычисляет $\mathrm{EXP}(X, b) = B_x$, получает те же ключи первой фазы, проверяет MAC и расшифровывает $CM$. Затем генерирует $(y, Y)$ и объединяет оба обмена, $\mathrm{EXP}(X, y)$ и $\mathrm{EXP}(X, b)$, в $\mathit{key\_seed}$. С помощью SHAKE-256 выводятся ключ для шифрования ответа $SM$ и поток ключей цепи. Реле отвечает $Y$ (32), $\mathit{AUTH}$ (32) и зашифрованным $SM$.
  \item Клиент проверяет $\mathit{AUTH}$, расшифровывает $SM$ и получает тот же поток ключей.
\end{enumerate}

\begin{attenzione}[Сообщение клиента не имеет прямой секретности]
$CM$ шифруется ключами, зависящими только от $B_x$, который выводится с использованием ключа onion реле. Тот, кто запишет обмен сообщениями и позже получит $b$, сможет расшифровать $CM$. Спецификация говорит об этом прямо: прямой секретностью обладают только данные реле. Поэтому $CM$ должен содержать только параметры согласования, но никогда не пользовательские данные.
\end{attenzione}

\subsection{Расширения}

Сообщения $CM$ и $SM$ используют общий формат~\spec{tor-spec}{Sending extensions in the circuit extension handshake}: один байт задаёт число расширений, затем для каждого идут тип (1 байт), длина (1 байт) и данные. Для создания цепи определены:
\begin{center}\small
\begin{tabularx}{\textwidth}{@{}cl>{\raggedright\arraybackslash}X@{}}
\toprule
Тип & Направление & Значение \\
\midrule
1 & клиент $\to$ реле & \file{CC\_FIELD\_REQUEST}: клиент запрашивает управление перегрузкой из \proposta[Предложения]{324}; пустое тело \\
2 & реле $\to$ клиент & \file{CC\_FIELD\_RESPONSE}: реле принимает; 1 байт, \file{sendme\_inc} \\
3 & клиент $\to$ реле & запрос субпротокола, например \texttt{[02 06]} для CGO (\texttt{Relay=6}) \\
\bottomrule
\end{tabularx}
\end{center}

\subsection{Какой протокол согласования ключей использует tor сегодня?}

В tor 0.4.9.13 тип согласования выбирается отдельно для каждого перехода (\file{circuit\_pick\_create\_handshake}). ntor-v3 используется, когда реле поддерживает управление перегрузкой и эта возможность включена, как обычно происходит в современной сети, либо когда требуется согласовать CGO. Иначе tor возвращается к ntor. На практике, как отмечает спецификация, современные клиенты почти всегда используют ntor-v3, а ntor сохраняется как вариант совместимости.

% ---------------------------------------------------------------------------
\section{От согласованных секретов к ключам цепи}\label{sec:circ-chiavi}

Каждое согласование ключей создаёт поток псевдослучайных байтов, из которого извлекаются ключи традиционной схемы шифрования ячеек RELAY, называемой \enquote{tor1}~\spec{tor-spec}{Setting circuit keys}:

\begin{center}
\setlength{\tabcolsep}{4pt}
\begin{tabular}{llll}
\toprule
Ключ & Байты & Назначение & Направление \\
\midrule
Df & 20 & начальное значение накопительного дайджеста (SHA-1) & клиент $\to$ реле \\
Db & 20 & начальное значение накопительного дайджеста (SHA-1) & реле $\to$ клиент \\
Kf & 16 & ключ AES-128-CTR & клиент $\to$ реле \\
Kb & 16 & ключ AES-128-CTR & реле $\to$ клиент \\
\bottomrule
\end{tabular}
\end{center}

Всего 72 байта (константа \file{CPATH\_KEY\_MATERIAL\_LEN} в коде tor). Способ формирования потока зависит от типа согласования ключей:
\begin{description}
  \item[KDF-TOR] (CREATE\_FAST и ранее TAP) $K = \mathrm{SHA1}(K_0 \| \texttt{00}) \,\|\, \mathrm{SHA1}(K_0 \| \texttt{01}) \,\|\, \ldots$. Первые 20 байт образуют $KH$, затем идут Df, Db, Kf и Kb. Однобайтовый счётчик ограничивает выход 5120 байтами, что значительно больше необходимого.
  \item[KDF-RFC5869] (ntor) использует HKDF с SHA-256 и \file{KEY\_SEED} как ключом; Df, Db, Kf и Kb занимают первые 72 байта.
  \item[SHAKE-256] (ntor-v3) использует первые 32 байта для шифрования $SM$; далее идут Df, Db, Kf и Kb.
\end{description}

Разные ключи в двух направлениях не избыточны. Если бы клиент и реле применяли один AES-CTR ключ в обе стороны, зашифрованные потоки повторно использовали бы одну и ту же гамму шифрования, и XOR двух ячеек раскрыл бы XOR двух открытых текстов.

CGO меняет схему ключей; это рассматривается в \cref{cap:10-cifratura-relay}.

% ---------------------------------------------------------------------------
\section{Лабораторная работа: построение реальных цепей}\label{sec:lab-circuiti}

\begin{laboratorio}[3]{Согласование ключей цепи}
\textbf{Цель.} Написать \file{labs/ru/код/torlab/ntor.py} с CREATE\_FAST, ntor и ntor-v3; проверить по официальным тестовым векторам; строить одношаговые цепи к реле лабораторной сети.\\
\textbf{Предварительные условия.} Лабораторные работы~1 и~2; работающая лабораторная сеть; библиотека \texttt{cryptography} (X25519, AES) и стандартный модуль \texttt{hashlib} (SHA-3, SHAKE).\\
\textbf{Оценочное время.} 3--4 часа.
\end{laboratorio}

\paragraph{Шаг 1: ntor-v3 в нескольких строках.} Реализация почти буквально следует спецификации. Сначала вспомогательные функции:
\begin{lstlisting}[language=Python]
NTOR3_PROTOID = b"ntor3-curve25519-sha3_256-1"
def _t(name):  return NTOR3_PROTOID + b":" + name.encode()
def _encap(s): return struct.pack(">Q", len(s)) + s
def _h3(s, t): return hashlib.sha3_256(_encap(t) + s).digest()
def _mac3(k, msg, t):
    return hashlib.sha3_256(_encap(t) + _encap(k) + msg).digest()
def _kdf3(s, t, n): return hashlib.shake_256(_encap(t) + s).digest(n)
\end{lstlisting}
Затем клиентская часть протокола согласования ключей:
\begin{lstlisting}[language=Python]
def ntor3_client_start(ID, B, message, ver=b"circuit extend", x=None):
    sk, X = x25519_generate(x)
    Bx = x25519_dh(sk, B)
    phase1 = _kdf3(Bx + ID + X + B + NTOR3_PROTOID + _encap(ver),
                  _t("kdf_phase1"), 64)
    enc_k1, mac_k1 = phase1[:32], phase1[32:]
    encrypted = aes_ctr(enc_k1, message)
    mac = _mac3(mac_k1, ID + B + X + encrypted, _t("msg_mac"))
    state = Ntor3State(sk, X, B, ID, Bx, mac, ver)
    return state, ID + B + X + encrypted + mac

def ntor3_client_finish(st, response, length=72):
    Y, auth, msg = response[:32], response[32:64], response[64:]
    secret_input = (x25519_dh(st.x, Y) + st.Bx + st.ID + st.B + st.X + Y
                    + NTOR3_PROTOID + _encap(st.ver))
    key_seed = _h3(secret_input, _t("key_seed"))
    verify = _h3(secret_input, _t("verify"))
    auth_input = (verify + st.ID + st.B + Y + st.X + st.mac
                  + _encap(msg) + NTOR3_PROTOID + b"Server")
    if not hmac.compare_digest(_h3(auth_input, _t("auth_final")), auth):
        raise HandshakeError("неверный ntor-v3 AUTH")
    stream = _kdf3(key_seed, _t("kdf_final"), 32 + length)
    return aes_ctr(stream[:32], msg), stream[32:]
\end{lstlisting}
Сравнение $\mathit{AUTH}$ выполняется \file{hmac.compare\_digest}, который занимает одинаковое время независимо от позиции первого несовпавшего байта. Обычное сравнение могло бы по времени раскрывать число совпавших начальных байтов.

\paragraph{Шаг 2: тестовые векторы.} Криптографический протокол согласования ключей нельзя проверять лишь тем, что он \enquote{работает}: две реализации с одной и той же ошибкой прекрасно взаимодействовали бы между собой. Поэтому используются \textbf{тестовые векторы} спецификации: фиксированные входы, включая якобы \enquote{случайные} ключи, и ожидаемые побайтовые выходы. \proposta{332} содержит полный вектор ntor-v3, а \proposta{216} --- несколько векторов HKDF ntor. Файл \file{labs/ru/код/torlab/test/test\_ntor.py} воспроизводит их; например:
\begin{lstlisting}[language=Python]
def test_ntor3_vectors_proposal_332():
    st, hdata = ntor3_client_start(h(V["ID"]), h(V["B"]), h(V["CM"]),
                                    ver=h(V["VER"]), x=h(V["x"]))
    assert hdata.hex() == V["client"]
    sm, stream = ntor3_client_finish(st, h(V["server"]), length=256)
    assert sm == h(V["SM"])
    assert stream[:32].hex() == V["keystream_start"]
\end{lstlisting}
Другие тесты проверяют, что изменённый $\mathit{AUTH}$ или другая строка \texttt{VER} приводят к ошибке согласования ключей.

\paragraph{Шаг 3: ключи реле.} Для согласования ключей нужны Ed25519-идентичность реле и \texttt{ntor-onion-key}. Настоящий клиент получает их из консенсуса и микродескрипторов (\cref{cap:12-directory}); мы берём их из тех же источников, используя файл \file{cached-microdescs.new} лабораторного клиента, где для каждого реле есть запись:
\begin{lstlisting}[numbers=none]
onion-key
-----BEGIN RSA PUBLIC KEY-----
MIGJAoGBAJet1Oxp+cvg/vgryEmJuVv2s++WxQt5Id1BQ9UNwCOVFcYvEN9KVklM
...
-----END RSA PUBLIC KEY-----
ntor-onion-key gsi2OisnQE96UkjOVvBKtWDLfwRdYG9d4DScg2wz4Qs
id ed25519 Kz1dj/u6/g7WvhHrt9Ds0jVZjKC+KHo26hWViXAJ1Pg
\end{lstlisting}
Обратите внимание на старый TAP-ключ, всё ещё публикуемый, но уже не используемый, и две base64-строки, которые нам действительно нужны.

\paragraph{Шаг 4: против настоящего реле.} Открываем канал к \texttt{test003r} с помощью лабораторной работы~1 и отправляем \cella{CREATE2} для каждого типа согласования ключей:
\begin{lstlisting}[language=Python]
c = open_channel("127.0.0.1", 5103, expected_identity=ed)
st, hd = ntor3_client_start(ed, B, encode_extensions([]))
c.send(Cell(0x80000001, CellCommand.CREATE2,
              create2_body(HTYPE_NTOR_V3, hd)))
r = c.receive()
sm, stream = ntor3_client_finish(st, hdata_created2(r.body))
keys = CircuitKeys.from_stream(stream)
\end{lstlisting}
Повторив опыт с другими вариантами, получаем:
\begin{lstlisting}[numbers=none]
<- CircID=0x80000001 CREATED2 тело=509 байт
   ntor-v3 ok, сообщение реле: []
<- CircID=0x80000002 CREATED2 тело=509 байт
   ntor-v3 с CC_FIELD_REQUEST ok, ответ: [(2, b'\x1f')]
<- CircID=0x80000003 CREATED2 тело=509 байт
   ntor ok
<- CircID=0x80000004 CREATED_FAST тело=509 байт
   CREATE_FAST ok
<- CircID=0x80000005 DESTROY тело=509 байт     (изменён MAC)
   причина: 1
\end{lstlisting}
Все неизменённые протоколы согласования ключей завершаются успешно. Во втором случае реле принимает управление перегрузкой и отвечает расширением 2 с \file{sendme\_inc}, равным 31 (\cref{cap:11-stream}). В последнем случае мы намеренно меняем один бит MAC: реле отклоняет согласование ключей ячейкой \cella{DESTROY} с причиной 1, \texttt{PROTOCOL}. Здесь реальную причину можно сообщать, потому что ошибку породило это реле; оно не пересылает чужую \cella{DESTROY} (\cref{sec:comandi}).

\paragraph{Неожиданное наблюдение.} В \cref{cap:07-canali} мы увидели, что новый канал не получает \cella{PADDING}. Теперь по этому каналу открыты четыре цепи, но за двадцать секунд ожидания по-прежнему не приходит ни одной ячейки. Причина находится в коде tor (\file{src/core/or/relay.c}): реле считает канал \enquote{используемым}, а значит пригодным для заполнения (padding), только если по нему идёт цепь, продолжающаяся за пределы этого реле. Наши одношаговые цепи не учитываются. В \cref{cap:10-cifratura-relay}, после расширения цепи появится заполнение.

% ---------------------------------------------------------------------------
\section*{Итоги}
\begin{itemize}
  \item Цепь строится по одному переходу: \cella{CREATE2}/\cella{CREATED2} с первым реле, затем \texttt{EXTEND2}/\texttt{EXTENDED2} через цепь для последующих. Каждое согласование ключей связывает клиента только с одним реле.
  \item \texttt{EXTEND2} содержит спецификаторы соединения (адрес, RSA- и Ed25519-идентичности); расширяющее реле обязано проверять эти идентичности.
  \item Протокол согласования ключей цепи аутентифицирует реле, но не клиента, и должен обеспечивать секретность и прямую секретность.
  \item TAP (1024-битные RSA и DH) удалён в tor 0.4.9.1-alpha. \cella{CREATE\_FAST} зависит от TLS и должен использоваться только для одношаговых цепей.
  \item ntor сочетает два обмена Curve25519: эфемерный ключ с эфемерным для прямой секретности и эфемерный ключ с ключом onion для аутентификации.
  \item ntor-v3 использует Ed25519-идентичность и семейство SHA-3, добавляя зашифрованные сообщения для согласования параметров вроде управления перегрузкой и CGO; сообщение клиента не имеет прямой секретности.
  \item Поток ключей даёт Df и Db (по 20 байт) и Kf и Kb (по 16 байт): ключ и дайджест для каждого направления.
\end{itemize}

% ---------------------------------------------------------------------------
\section*{Упражнения}

\begin{esercizio}[1]\label{es:09-dimensioni}
Вычислите длину HDATA в \cella{CREATE2} для ntor и для ntor-v3 с трёхбайтовым сообщением. Каков максимальный размер сообщения $CM$ ntor-v3 внутри \cella{CREATE2}?
\end{esercizio}

\begin{esercizio}[1]\label{es:09-fast}
Почему \cella{CREATE\_FAST} нельзя использовать для второго перехода цепи? Что могло бы сделать первое реле?
\end{esercizio}

\begin{esercizio}[1]\label{es:09-direzioni}
Почему клиент и реле используют разные AES-ключи в двух направлениях (Kf и Kb)? Опишите, что произошло бы при использовании одного ключа.
\end{esercizio}

\begin{esercizio}[2]\label{es:09-extend2}
Запишите в шестнадцатеричном виде заголовок тела сообщения \texttt{EXTEND2} --- число спецификаторов и спецификатор IPv4 --- для расширения цепи до лабораторного реле \texttt{test004r}. Сколько байт занимает полное сообщение с ntor-v3 и однобайтовым $CM$?
\end{esercizio}

\begin{esercizio}[2]\label{es:09-ds}
Предположим, в ntor клиент использует только $\mathrm{EXP}(Y, x)$ и опускает $\mathrm{EXP}(B, x)$. Опишите конкретную атаку промежуточного реле, пересылающего согласование ключей следующему переходу. Теперь предположим, что клиент вместо этого опускает $\mathrm{EXP}(Y, x)$: какое свойство теряется?
\end{esercizio}

\begin{esercizio}[2]\label{es:09-ver}
Каково назначение строки \texttt{VER} в ntor-v3? Представьте, что другой протокол, не Tor, использует то же согласование ключей и те же ключи: что могло бы произойти без \texttt{VER}?
\end{esercizio}

\begin{esercizio}[2]\label{es:09-cc}
Повторите шаг 4 лабораторной работы с запросом субпротокола \texttt{[02 06]} (CGO) вместе с \file{CC\_FIELD\_REQUEST}. Принимает ли его реле? С помощью спецификации \proposta{359} объясните, почему CGO также требует управления перегрузкой.
\end{esercizio}

\begin{esercizio}[3]\label{es:09-ntor-ref}
Репозиторий tor содержит эталонную реализацию ntor на Python (\file{src/test/ntor\_ref.py}). Используйте её для создания полных тестовых векторов протокола ntor для клиента и реле и добавьте их в \file{test\_ntor.py}.
\end{esercizio}

\begin{esercizio}[3]\label{es:09-relay}
Реализуйте \emph{сторону реле} ntor-v3 (функции \file{server\_part1} и \file{server\_part2}) и проверьте её по векторам \proposta[Предложения]{332}. Затем выполните полный обмен сообщениями между своим клиентом и реле в памяти.
\end{esercizio}

% ---------------------------------------------------------------------------
\section*{Дополнительная литература}
\begin{itemize}
  \item \titolo{Tor Protocol Specification}, главы \enquote{Creating and extending circuits} и \enquote{Setting circuit keys}~\cite{torspec}.
  \item I.~Goldberg, D.~Stebila, B.~Ustaoglu, \titolo{Anonymity and one-way authentication in key exchange protocols}~\cite{goldberg2013ntor}: статья, определяющая ntor и доказывающая его безопасность.
  \item Предложения 216 (ntor), 332 (ntor-v3) и 350 (удаление TAP).
\end{itemize}
